\section*{Capítulo \ref{ch:b3:cell-cycle-apoptosis} --- Control del ciclo celular y muerte celular programada}

\begin{solution}{exo:b3:cell-cycle-apoptosis:1}
G1, S, G2, M. Avance de G1 y \omterm{def:b3:cell-cycle-apoptosis:checkpoints}{punto de restricción}: \omterm{def:b3:cell-cycle-apoptosis:cdk}{ciclina} D--Cdk4/6;
G1/S: \omterm{def:b3:cell-cycle-apoptosis:cdk}{ciclina} E--Cdk2; S: \omterm{def:b3:cell-cycle-apoptosis:cdk}{ciclina} A--Cdk2; G2/M: \omterm{def:b3:cell-cycle-apoptosis:cdk}{ciclina} B--Cdk1. El
\omterm{def:b3:cell-cycle-apoptosis:cdk}{complejo promotor de la anafase} (APC/C) destruye la \omterm{def:b3:cell-cycle-apoptosis:cdk}{ciclina} B y la
securina y termina la mitosis.
\end{solution}

\begin{solution}{exo:b3:cell-cycle-apoptosis:2}
Hartwell: los mutantes \emph{cdc} de la levadura de gemación, detenidos
en etapas definidas, y el concepto de Start. Nurse: \emph{cdc2} de la
levadura de fisión como la quinasa que marca el momento de la mitosis, y
su homólogo humano (Cdk1) rescatando a la levadura --- la conservación.
Hunt: la \omterm{def:b3:cell-cycle-apoptosis:cdk}{ciclina} en huevos de erizo de mar, una proteína sintetizada a lo
largo del ciclo y destruida en cada división.
\end{solution}

\begin{solution}{exo:b3:cell-cycle-apoptosis:3}
\omterm{def:b3:cell-cycle-apoptosis:checkpoints}{Punto de restricción} (G1 tardía): factores de crecimiento, nutrientes,
tamaño; actúa sobre la \omterm{def:b3:cell-cycle-apoptosis:cdk}{ciclina} D--Cdk4/6 y sobre Rb. G2/M: ADN intacto y
replicado; \omterm{def:b3:dna-repair:ddr}{ATM}/ATR inhiben a Cdc25 y mantienen a Cdk1 fosforilada e
inactiva. Ensamblaje del huso: todos los cinetocoros unidos; los
cinetocoros sin unir forman complejos Mad2--Cdc20 que inhiben el APC/C.
\end{solution}

\begin{solution}{exo:b3:cell-cycle-apoptosis:4}
\omterm{def:b3:cell-cycle-apoptosis:apoptosis}{Apoptosis}: la célula se encoge, la \omterm{def:b3:chromatin-epigenetics:nucleosome}{cromatina} se condensa, el ADN se corta
en escalera, la membrana forma ampollas pero sigue sellada, se expone
\omterm{def:b3:cell-cycle-apoptosis:apoptosis}{fosfatidilserina}, el cadáver es engullido y no hay inflamación. \omterm{def:b3:cell-cycle-apoptosis:apoptosis}{Necrosis}:
la célula se hincha, la membrana se rompe, el contenido se derrama, el
ADN se degrada al azar y hay inflamación. Ejecutoras: las \omterm{def:b3:cell-cycle-apoptosis:apoptosis}{caspasas},
cisteína-proteasas que cortan detrás de un aspartato.
\end{solution}

\begin{solution}{exo:b3:cell-cycle-apoptosis:5}
Subida de $5$ a $30$ a $1$ por minuto: \qty{25}{min}. Degradación de $30$
a $5$ con semivida \qty{2}{min}: $\log_{2}(30/5)\times 2 =
\qty{5.2}{min}$. Período de unos \qty{30}{min}, con la mitosis en el
\qty{17}{\%}.
\end{solution}

\begin{solution}{exo:b3:cell-cycle-apoptosis:6}
(a) $C$ nunca puede bajar de $C_{1}$: queda bloqueada en mitosis, en la
rama alta. (b) Sin fosforilación inhibidora: el umbral $C_{2}$ baja y la
mitosis empieza antes, con un tamaño pequeño (el fenotipo ``wee'').
(c) Cdk1 permanece fosforilada: la rama alta es inalcanzable, y hay
parada en G2 con células alargadas. (d) Como en (b), y además el \omterm{def:b3:cell-cycle-apoptosis:checkpoints}{punto de
control} G2/M, que actúa a través de Wee1/Cdc25, ya no puede retener a la
célula.
\end{solution}

\begin{solution}{exo:b3:cell-cycle-apoptosis:7}
Núcleo en G1 en citoplasma de fase S: sus orígenes están licenciados y el
citoplasma aporta \omterm{def:b3:cell-cycle-apoptosis:cdk}{ciclina} E/A--Cdk2 activa, de modo que entra en S de
inmediato. Núcleo en G2 en citoplasma de fase S: sus orígenes ya
dispararon y no se relicenciaron (los factores de licencia los destruye o
los inhibe la CDK), así que no puede replicar de nuevo; espera hasta que
la pareja llega a G2 y ambos entran juntos en mitosis.
\end{solution}

\begin{solution}{exo:b3:cell-cycle-apoptosis:8}
El cinetocoro sin unir es un catalizador: convierte de continuo la Mad2 a
su conformación activa y produce un inhibidor difusible de Cdc20 que se
extiende por la célula y mantiene apagado el APC/C en todas partes ---
con un solo cinetocoro basta. Sin Mad2 el APC/C se activa en cuanto Cdk1
lo ha encendido, estén o no unidos los cromosomas: la anafase empieza con
cromosomas sin unir, las hijas son aneuploides y el organismo (o la línea
celular) muere por pérdida de cromosomas.
\end{solution}

\begin{solution}{exo:b3:cell-cycle-apoptosis:9}
Sin \emph{ced-9}: mueren células que deberían vivir --- letalidad
embrionaria por exceso de muerte. Sin \emph{ced-3}: no ocurre ninguna de
las $131$ muertes, las células sobreviven y se diferencian, y el gusano
es viable. Sin ambos: no hay muerte --- el fenotipo de \emph{ced-3}. Como
quitar la ejecutora anula el efecto de quitar al guardián, CED-9 actúa
aguas arriba, manteniendo inactivos a CED-4 y CED-3; cuando desaparece,
estos quedan activos, salvo que también falten.
\end{solution}

\begin{solution}{exo:b3:cell-cycle-apoptosis:10}
Con $A = aC$, $\mathrm{d}C/\mathrm{d}t = k_{s} - k_{d}aC^{2}$, una
ecuación de una sola variable con estado estacionario
$C^{*} = \sqrt{k_{s}/(k_{d}a)}$ y pendiente $-2k_{d}aC^{*} < 0$ en él:
toda desviación decae de forma monótona, y una sola variable de primer
orden no puede oscilar. La curva en forma de S da dos ramas estables
separadas por una zona prohibida, de modo que el estado tiene que saltar
y no puede asentarse; un retraso explícito en la retroalimentación daría
oscilaciones por otra vía --- el freno llegando demasiado tarde, el
exceso, y así sucesivamente --- como en muchos ritmos hormonales.
\end{solution}

\begin{solution}{exo:b3:cell-cycle-apoptosis:11}
Cada sensor que llega lo captura un guardián, de modo que no ocurre nada
hasta que los $N$ guardianes están ocupados, hacia $t \approx N/r$; los
sensores siguientes quedan libres, activan a Bax/Bak, cuya formación de
poros sube abruptamente con su número, y las mitocondrias se abren en
minutos. Un estímulo más fuerte acorta el tiempo solo a través de $r$;
por encima del umbral el desenlace es el mismo. El venetoclax ocupa
guardianes y reduce el $N$ efectivo: el tiempo se acorta, y una célula
cebada --- ya cerca de $N$ --- muere de inmediato.
\end{solution}

\begin{solution}{exo:b3:cell-cycle-apoptosis:12}
El exceso de Bcl-2 bloquea la \omterm{def:b3:cell-cycle-apoptosis:pathways}{vía intrínseca}: los linfocitos~B que
deberían morir cuando cesan sus señales de crecimiento sobreviven, y la
población crece por fallo de la muerte, al ritmo lento con que se
fabrican esas células --- indolente, hasta que una segunda lesión añade
proliferación. La pérdida de Rb elimina el \omterm{def:b3:cell-cycle-apoptosis:checkpoints}{punto de restricción}: los
precursores retinianos recorren el ciclo sin factores de crecimiento y se
dividen tan deprisa como permite el motor --- agresivo. Los dos tumores
son los dos programas fallando cada uno por su cuenta: un control de
muerte y un control de división, y basta con perder cualquiera de ellos
para tener un tumor, más rápido con el segundo.
\end{solution}

\begin{solution}{pb:b3:cell-cycle-apoptosis:1}
\textbf{1.} $10^{7}\times 3500/5 = 7\times 10^{9}$ células al día, unas
\num{80000} por segundo.
\textbf{2.} $3500/5 = 700$ por cripta y día; $22\times 2^{5} = 704$.
\textbf{3.} $80\times 365 \approx \num{29000}$ divisiones; $\times 0.6
\approx \num{17500}$ mutaciones.
\textbf{4.} Las células de tránsito se desprenden en pocos días y se
llevan sus mutaciones consigo; la célula madre persiste toda la vida, de
modo que se conserva cada una de sus mutaciones y cada división añade más
--- ciclar despacio las minimiza.
\textbf{5.} La vellosidad se vacía en los $5$ días de su tránsito normal;
reconstruirla cuesta un día de recuperación, cinco divisiones de
\qty{12}{h} y la migración vellosidad arriba --- de cuatro a cinco días,
durante los cuales la barrera está desnuda: las rápidas divisiones de
tránsito del intestino lo convierten en una víctima temprana de la
radiación.
\textbf{6.} Un cadáver apoptótico sellado no libera ningún contenido a
una luz llena de bacterias y no provoca inflamación; la \omterm{def:b3:cell-cycle-apoptosis:apoptosis}{necrosis}
derramaría enzimas y señales e inflamaría la mucosa cada vez que se
desprende una célula.
\textbf{7.} $(40 - 8)/2 = \qty{16}{h}$.
\textbf{8.} $\log_{2}(40/8)\times \qty{10}{min} = 2.3\times 10 =
\qty{23}{min}$.
\textbf{9.} Período de unas \qty{16.4}{h}; Cdk1 activa el \qty{2.3}{\%}
del tiempo.
\textbf{10.} No $k_{s}$ por sí solo, sino el tiempo retenido en los
umbrales: la célula somática espera en el \omterm{def:b3:cell-cycle-apoptosis:checkpoints}{punto de restricción} a los
factores de crecimiento y en G2/M a su ADN, y la célula madre espera más
que la de tránsito. El huevo de rana no tiene \omterm{def:b3:cell-cycle-apoptosis:checkpoints}{puntos de control} y tiene
un citoplasma abastecido de todo, de modo que solo el oscilador desnudo
fija su período.
\textbf{11.} En la rama baja, a un nivel de \omterm{def:b3:cell-cycle-apoptosis:cdk}{ciclina} por encima del
$C_{2}$ normal: el \omterm{def:b3:cell-cycle-apoptosis:checkpoints}{punto de control} ha desplazado la curva en S hacia la
derecha al inhibir a Cdc25, de modo que el salto no ocurre. Al repararse,
Cdc25 queda liberada, el umbral cae por debajo de la \omterm{def:b3:cell-cycle-apoptosis:cdk}{ciclina} acumulada y
la célula entra de inmediato en mitosis --- que es por lo que las células
liberadas de un \omterm{def:b3:cell-cycle-apoptosis:checkpoints}{punto de control} entran en mitosis de forma sincrónica.
\textbf{12.} No se degradan ni la \omterm{def:b3:cell-cycle-apoptosis:cdk}{ciclina} B ni la securina: la célula se
queda en metafase con la cohesina intacta y Cdk1 activa, indefinidamente;
muere ahí o acaba escapándose con el genoma sin segregar.
\textbf{13.} $10^{11}$ divisiones al día.
\textbf{14.} Con el \omterm{def:b3:cell-cycle-apoptosis:checkpoints}{punto de control},
$10^{11}/5\times 10^{4} = 2\times 10^{6}$ hijas aneuploides al día; sin
él, $2\times 10^{9}$.
\textbf{15.} $2\times 10^{4}$ al día; a lo largo de $80$ años,
$6\times 10^{8}$.
\textbf{16.} $10^{9}/50 = 2\times 10^{7}$ malas segregaciones al día:
cada día el tumor prueba veinte millones de cariotipos nuevos, y la
selección conserva los que crecen más deprisa o resisten un fármaco.
\textbf{17.} Si toda división segrega mal, ningún linaje celular
embrionario conserva un genoma euploide y el embrión muere. La pérdida
parcial da una aneuploidía ocasional en células que la sobreviven (sobre
todo sin p53), una inestabilidad cromosómica que aporta variación a un
tumor sin matar al organismo.
\textbf{18.} Un error meiótico mete el cromosoma de más en el gameto, de
modo que el cigoto y todas las células derivadas de él son trisómicos; un
error mitótico ocurre en una sola célula somática y solo lo hereda su
clon.
\textbf{19.} $10^{11}\times \qty{1}{ng} = \qty{100}{g}$ al día,
\qty{36}{kg} al año --- alrededor de la mitad de la masa corporal cada
año.
\textbf{20.} \qty{20}{g} de proteína al día, cerca de un tercio de la
ingesta alimentaria y un pequeño porcentaje del recambio proteico total
del cuerpo.
\textbf{21.} $10^{11}$ cadáveres a $24$ por macrófago y día: $4\times
10^{9}$ macrófagos a tiempo completo, el \qty{4}{\%} del tiempo de los
$10^{11}$ macrófagos.
\textbf{22.} Un cadáver lisado libera proteasas, ADN, ATP y otras señales
de alarma que causan inflamación y pueden provocar autoinmunidad contra
antígenos nucleares; el cadáver íntegro se anuncia con \omterm{def:b3:cell-cycle-apoptosis:apoptosis}{fosfatidilserina}
en su monocapa externa (y perdiendo sus señales de ``no me comas'') y
atrae a los fagocitos con nucleótidos liberados.
\textbf{23.} La \omterm{def:b3:cell-cycle-apoptosis:pathways}{vía intrínseca} queda bloqueada en la mitocondria. La ruta
extrínseca sigue funcionando donde la caspasa-8 activa directamente a la
caspasa-3, aunque se pierde su amplificación a través de Bid. Las células
se acumulan por no morir más que por dividirse más deprisa ---
crecimiento lento --- hasta que una lesión posterior (a menudo
\emph{MYC}) añade proliferación.
\textbf{24.} La \omterm{def:b3:cell-cycle-apoptosis:apoptosis}{apoptosis} tarda horas, necesita ATP y pasa por etapas que
pueden bloquearse --- inhibidores de \omterm{def:b3:cell-cycle-apoptosis:apoptosis}{caspasas}, el umbral mitocondrial
---, de modo que las neuronas de la penumbra pueden rescatarse si se
tratan a tiempo; las del núcleo mueren por fallo energético en minutos,
sin programa que interrumpir.
\textbf{25.} Período del oscilador desnudo, unas \qty{16}{h}; unas $2\times
10^{4}$ células aneuploides en división al día en un cuerpo normal; y
\qty{100}{g} de células recicladas por \omterm{def:b3:cell-cycle-apoptosis:apoptosis}{apoptosis} cada día.
\end{solution}
